\section{Saturation of determinant support and extension across the boundary}

\subsection*{M-020. Two precise statements that do not compare the original ranks}
There is a known saturation theorem for the actual determinant closure,
as well as a known localization theorem relating its orbit and closure.
Their short consequences below specify what they give in the report's
notation. They must be distinguished from ordinary Kronecker saturation,
and from a comparison of permanent and determinant multiplicities.
No historical priority is asserted for these elementary consequences.

\paragraph{Known support theorem and a finite uniform-stretch reduction.}
Let
\[
 S_n=\{\lambda:K_n(\lambda)>0\},\qquad
 T_n=\{\lambda:\ell(\lambda)\leq n,\ n\mid |\lambda|\}.
\]
Zero is included in both monoids. B\"urgisser--Ikenmeyer--H\"uttenhain,
\emph{Permanent versus determinant: not via saturations},
arXiv:1501.05528v2, Theorem~1, printed p.~3, proves
\[
                         T_n\subseteq\operatorname{Sat}(S_n)
                         \qquad(n>2).                         \tag{M-020.1}
\]
Here saturation is taken in the group generated by the monoid,
as defined immediately before that theorem. In particular each
\(\lambda\in T_n\) has some positive integral stretch in \(S_n\).
The primary source is \url{https://arxiv.org/pdf/1501.05528v2}.
This is a cited theorem, not a result derived here from Burnside
effectivity or K-theoretic positivity.

For clarity, there is even a uniform stretching integer for each
fixed \(n\):
\[
 \text{there exists }c_n\geq1\text{ such that }
            K_n(c_n\lambda)>0\quad\text{for every }\lambda\in T_n.
                                                               \tag{M-020.2}
\]
One finite reduction for choosing it uses
\[
 H_n=\{\theta\ne0:\ell(\theta)\leq n,\ \theta_1\leq n,
                                      \ n\mid |\theta|\}.
                                                               \tag{M-020.3}
\]
For each \(\theta\in H_n\), choose any \(c_\theta\geq1\)
with \(K_n(c_\theta\theta)>0\), which exists by (M-020.1).
Then \(c_n=\operatorname{lcm}_{\theta\in H_n}c_\theta\)
works. This specifies a finite gate; it supplies no numerical
bound on the choices \(c_\theta\).

\begin{proof}[The finite reduction]
Regard a partition in \(T_n\) as a multiset of columns with
heights in \(\{1,\ldots,n\}\). Any sequence of at least \(n\)
columns contains a nonempty consecutive subsequence of length
at most \(n\) whose total height is divisible by \(n\): among
the first \(n+1\) partial sums, including zero, two agree
modulo \(n\). Remove that subsequence. Its column heights,
sorted, define an element of \(H_n\); removal preserves
divisibility of the remaining total height. Continue until
fewer than \(n\) columns remain, which themselves give an
element of \(H_n\), unless the remainder is empty.
Thus every element of \(T_n\) is a sum of elements of \(H_n\).
The determinant closure is irreducible. Its coordinate ring
is therefore a domain, and products of nonzero highest-weight
vectors show that \(S_n\) is closed under addition. Multiples
of \(c_\theta\theta\) and then sums of those multiples prove
(M-020.2) for the stated least common multiple.
\end{proof}

For internal padding, \(q^m_\lambda(d[n])>0\) implies
\(\ell(\lambda)\leq m^2\) and \(|\lambda|=nd\).
Consequently, when \(n\geq m^2\) and \(n>2\), every possible
permanent type already belongs to \(\operatorname{Sat}(S_n)\).
Replacing determinant support by its saturation therefore erases
any occurrence distinction in this range. This statement gives
neither \(K_n(\lambda)>0\) at the original weight nor an inequality
between two positive multiplicities. It also does not say that
every sufficiently large integer stretch is positive: a uniform
multiple and all sufficiently large integers are different assertions.
In the stronger BIP range the separately cited no-occurrence theorem
does give original-weight positivity; the positive-rank comparison
still remains.

\paragraph{Known boundary localization and its all-copy consequence.}
Put \(R=n^2\), \(G=\operatorname{GL}_R\),
\(X=\overline{G\cdot\det_n}\), and \(O=G\cdot\det_n\).
Write \(A=\mathbb C[X]\) and \(B=\mathbb C[O]\).
B\"urgisser--Ikenmeyer, \emph{Fundamental invariants of orbit closures},
arXiv:1511.02927v2, Definition~3.8 and Proposition~3.9,
printed pp.~9--10, supply a homogeneous nonzero fundamental
\(\operatorname{SL}_R\)-invariant \(\Phi\in A\) such that
\[
             O=X\setminus V(\Phi),\qquad B=A[\Phi^{-1}].
                                                               \tag{M-020.4}
\]
The source's determinant polystability is Corollary~2.9,
printed pp.~6--7. The primary PDF is
\url{https://arxiv.org/pdf/1511.02927v2}.
Write \(e=\deg\Phi\). In the polynomial coordinate convention
used here, \(\Phi\) spans a one-dimensional representation
of highest weight \(\delta\mathbf1_R\), where
\[
                 \delta\in\mathbb Z_{>0},\qquad e=n\delta.
                                                               \tag{M-020.5}
\]
Indeed an \(\operatorname{SL}_R\)-invariant polynomial
\(\operatorname{GL}_R\)-character is a nonnegative determinant
power, and comparison with the scalar action gives
\(R\delta=ne\). No explicit value of \(e\) is assumed.

Fix \(\lambda\vdash nd\), and let
\(b_\lambda=[\mathbb S_\lambda\mathbb C^R:B]\), which is
the symmetric Kronecker orbit multiplicity in this report.
Then
\[
 K_n(\lambda+k\delta\mathbf1_R)
 \ \text{is nondecreasing in }k\geq0,\quad
 K_n(\lambda+k\delta\mathbf1_R)\leq b_\lambda,
                                                               \tag{M-020.6}
\]
and, for some \(k_0=k_0(n,\lambda,\Phi)\),
\[
             K_n(\lambda+k\delta\mathbf1_R)=b_\lambda
                         \qquad\text{for every }k\geq k_0.
                                                               \tag{M-020.7}
\]
The outer degree on the left is \(d+ke\). This is a
rectangular translation of the weight, not its dilation.

\begin{proof}[The complete multiplicity-space statement]
Let \(A^U_\mu\) and \(B^U_\mu\) be the highest-weight spaces
for a fixed maximal unipotent subgroup \(U\). Multiplication
by \(\Phi\) is injective because \(A\) is a domain, and shifts
the weight by \(\delta\mathbf1_R\). Division by \(\Phi^k\)
is defined in \(B\) and identifies the image of
\(A^U_{\lambda+k\delta\mathbf1_R}\) with a subspace of
\(B^U_\lambda\). These subspaces are nested as \(k\) grows.
They have dimensions equal to the left side of (M-020.6).

The orbit highest-weight space is finite-dimensional; in this
case its dimension is the known symmetric Kronecker coefficient.
Choose a basis \(f_1,\ldots,f_{b_\lambda}\). By (M-020.4),
for each basis vector some \(\Phi^{k_i}f_i\) belongs to \(A\).
It remains a highest-weight vector, of the shifted weight,
because \(\Phi\) is a semi-invariant. Taking
\(k_0=\max_i k_i\) makes every basis vector occur in the
divided image for every \(k\geq k_0\). The upper bound in
(M-020.6) then gives equality. When \(b_\lambda=0\), take
\(k_0=0\); both sides are zero for all \(k\).
\end{proof}

For example, F-025's type has orbit multiplicity ten and
closure multiplicity five. It follows that its translated
closure multiplicities eventually equal ten, although no
value of the first successful translation is supplied here.
This argument does not determine the dilates in A-030.

There is a direct padding limitation. If \(m<n\) and \(k>0\),
then \(\lambda+k\delta\mathbf1_{n^2}\) has exactly \(n^2\)
nonzero parts. Thus, at its corresponding degree,
\[
 q^m_{\lambda+k\delta\mathbf1_{n^2}}((d+ke)[n])=0,
                                                               \tag{M-020.8}
\]
by the permanent's \(m^2\) essential-variable bound. Hence
clearing determinant-boundary denominators in this manner
leaves the possible internally padded permanent support.
The localization theorem does not compare the original two
positive ranks and is not a way to transfer an unpadded
obstruction to the BIP range.

\paragraph{Grothendieck dictionary.}
The support statement concerns the character grading on
\(A^U\); its saturation is a lattice-and-cone operation,
not the class of a virtual representation in \(K_0\).
The coordinate-comodule and weight-space dictionary is
SGA~3, I, 4.7.2--4.7.3, pp.~34--36.
The boundary statement is the distinguished affine open
\(D(\Phi)\subseteq\operatorname{Spec}A\), with coordinate
ring \(A_\Phi\); it uses ordinary localization and the
finite-dimensionality of each Schur multiplicity space.
The Schur operations themselves are those of SGA~6, VI,
Theorem~3.3, pp.~392--393. The proof explicitly clears
denominators of a basis of actual sections. It is therefore
a statement about extension of sections across the boundary,
with an unspecified extension threshold, rather than an
effectivity consequence of a virtual \(\lambda\)-identity.
